\chapter{Geometric Number Theory}

\section{Pick's Theorem}

\begin{definition}
  Consider two points $(a,b)$ and $(c,d)$ in the plane. The line segment $L((a,b),(c,d))$ is defined by
  \begin{equation*}
    L((a,b),(c,d))= \{t(a,b)+(1-t)(c,d)\mid 0\leq t\leq 1\}.
  \end{equation*}
\end{definition}

\begin{definition}
  A \emph{simple polygon with integer vertices} in the plane consists of an ordered list of points $(x_i,y_i)$  with integer coordinates $x_i$ and $y_i$, where $0\leq i < n$, such that the line segments $L_i:=L((x_i,y_i),(x_{i+1},y_{i+1}))$, satisfy the following condition:
  \begin{equation*}
    L_i\cap L_j=
    \begin{cases}
      \{(x_j,y_j)\} & \text{if $j-i\equiv 1\pmod{n}$}, \\
      \{(x_i,y_i)\} & \text{if $i-j\equiv 1\pmod{n}$}, \\
      \varnothing & \text{otherwise.}
    \end{cases}
  \end{equation*}
  In other words, a simple polygon with integer vertices in the plane consists of an ordered list of 
\end{definition}

\begin{theorem}[Pick's Theorem]
  If $P$ is a polygon whose vertices have integer coordinates, then
  
\end{theorem}

\section{Minkowski's Convex Body Theorem}

\begin{theorem}[Minkowski's Convex Body Theorem]
  Consider a lattice $\Lambda\subseteq \mathbb{R}^n$, and consider a convex, centrally symmetric, measurable subset $K\subseteq \mathbb{R}^n$ with
  \begin{equation*}
    \mathrm{vol}(K)\geq 2^n\mathrm{det}(\Lambda).
  \end{equation*}
  Then $K$ contains a nonzero element of $\Lambda$.
\end{theorem}

\section{Blichfeldt's Principle}

\begin{theorem}[Blichfeldt's Principle]
  If $\mathrm{vol}(S)\geq\mathrm{det}(\Lambda)$, then there exist distinct $x,y\in S$ such that $x-y\in\Lambda$.
\end{theorem}

\section{Erhardt's Theorem}

%%% Local Variables:
%%% mode: latex
%%% TeX-master: "number-theory"
%%% TeX-engine: luatex
%%% End:
